% Target after editing: approximately 0.75--1 page including tables and reproduction protocol.
\section{Experiments}

\subsection{Experimental Setup}
\pending{Report the six datasets, tensor shapes, baselines, normalization, input windows, forecast horizons, evaluation metrics, hardware, and any tuning settings actually used in the selected discussion. Cite exact table/section locations from the paper.}

\begin{table}[htbp]
  \centering
  \caption{Selected Forecasting Results Supported by the Paper Only}
  \label{tab:paper-reference-format}
  \small
  \begin{tabular}{@{}lccc@{}}
    \toprule
    Dataset & Horizon & MAE & RMSE \\
    \midrule
    \pending{Choose one} & \pending{value} & \pending{value} & \pending{value} \\
    \bottomrule
  \end{tabular}
\end{table}

\subsection{Results from the Original Paper}
\pending{Summarize forecasting accuracy, ablation, scalability, and effectiveness evidence. Use values supported by the paper and mark them explicitly as ``paper-reported,'' not ours. Include comparisons against the relevant baselines.}

\subsection{Our Reproduction}
\pending{Document code/data versions, environment and hardware, preprocessing, train/validation/test indexing, window boundaries, stopping rules, seeds, and the exact evaluation loop.}

\pending{Insert measured MAE/RMSE, runtime and ablation comparisons from team experiment logs. If a result is unavailable or differs from the paper, explain the discrepancy rather than inventing a value or claiming success.}

\begin{table}[htbp]
  \centering
  \caption{Team Reproduction Results versus Paper References}
  \label{tab:team-reproduction-format}
  \small
  \begin{tabular}{@{}lcc@{}}
    \toprule
    Metric & Paper Reference & Our Run \\
    \midrule
    MAE & \pending{cited value} & \pending{measured value} \\
    RMSE & \pending{cited value} & \pending{measured value} \\
    Runtime & \pending{cited setup/value} & \pending{measured value} \\
    \bottomrule
  \end{tabular}
\end{table}
